% =========================================================================
\section{Raíces múltiples}\label{sec:multiples}
% =========================================================================

Si $\lambda_k$ tiene multiplicidad $m_k>1$, la exponencial $\eul^{\lambda_kx}$
está en $\Ker(D-\lambda_k)^{m_k}$, pero veremos que ese núcleo tiene
dimensión $m_k$: faltan $m_k-1$ soluciones.

% -------------------------------------------------------------------------
\subsection{Conjugación por la exponencial}\label{sec:conjugacion}
% -------------------------------------------------------------------------

\begin{lema}[«Derivada covariante»]\label{lem:covariante}
Para todo $\lambda\in\C$, todo $m\in\N$ y toda $g\in\mathcal C^{m}(I)$,
\begin{equation}\label{eq:covariante}
  (D-\lambda)\bigl(\eul^{\lambda x}g\bigr)=\eul^{\lambda x}g',
  \qquad
  (D-\lambda)^{m}\bigl(\eul^{\lambda x}g\bigr)=\eul^{\lambda x}g^{(m)} .
\end{equation}
Es decir, $\eul^{-\lambda x}(D-\lambda)\,\eul^{\lambda x}=D$: conjugado por la
exponencial, $D-\lambda$ actúa como la derivada ordinaria.
\end{lema}

\begin{proof}
$(D-\lambda)(\eul^{\lambda x}g)=\lambda\eul^{\lambda x}g+\eul^{\lambda x}g'-\lambda\eul^{\lambda x}g=\eul^{\lambda x}g'$.
La segunda igualdad se obtiene aplicando la primera $m$ veces (inducción
sobre $m$).
\end{proof}

\begin{proposicion}[Núcleo de $(D-\lambda)^{m}$]\label{prop:nucleo}
Para todo $\lambda\in\C$ y todo $m\ge1$,
\begin{equation}\label{eq:nucleo}
  \Ker(D-\lambda)^{m}=\spanop\bigl\{\eul^{\lambda x},\,x\,\eul^{\lambda x},\,\dots,\,x^{m-1}\eul^{\lambda x}\bigr\},
  \qquad \dim\Ker(D-\lambda)^{m}=m .
\end{equation}
\end{proposicion}

\begin{proof}
Toda función se escribe como $y=\eul^{\lambda x}g$ con $g=\eul^{-\lambda x}y$.
Por el \cref{lem:covariante},
\[
  (D-\lambda)^{m}y=0\iff g^{(m)}=0\iff g(x)=c_0+c_1x+\dots+c_{m-1}x^{m-1}.
\]
Las $m$ funciones $x^{j}\eul^{\lambda x}$ son independientes porque
$\eul^{\lambda x}\neq0$ y los monomios lo son.
\end{proof}

% -------------------------------------------------------------------------
\subsection{Solución general}\label{sec:general-multiples}
% -------------------------------------------------------------------------

\begin{lema}[Independencia]\label{lem:independencia}
Si $\lambda_1,\dots,\lambda_r\in\C$ son distintas y $P_1,\dots,P_r$ son
polinomios con $\sum_kP_k(x)\,\eul^{\lambda_kx}\equiv0$ en $I$, entonces
$P_1=\dots=P_r=0$.
\end{lema}

\begin{proof}
Por inducción sobre $r$; el caso $r=1$ es inmediato. Multiplicando por
$\eul^{-\lambda_rx}$ y derivando $\deg P_r+1$ veces se elimina el último
término. Como $D\bigl(P\eul^{\mu x}\bigr)=(P'+\mu P)\,\eul^{\mu x}$ y
$\mu=\lambda_k-\lambda_r\neq0$ para $k<r$, cada derivación transforma
$P_k\eul^{(\lambda_k-\lambda_r)x}$ en $Q_k\eul^{(\lambda_k-\lambda_r)x}$ con
$\deg Q_k=\deg P_k$. La hipótesis de inducción da $Q_k=0$, luego $P_k=0$
para $k<r$, y entonces $P_r\eul^{\lambda_rx}\equiv0$ obliga a $P_r=0$.
\end{proof}

\begin{teorema}[Base de $\Ker\Lop$ con coeficientes constantes]\label{teo:base}
Con la factorización \eqref{eq:factorizacion-L}, las $n$ funciones
$x^{j}\eul^{\lambda_kx}$, $0\le j<m_k$, $1\le k\le r$, forman una base de
$\Ker\Lop$, y
\begin{equation}\label{eq:yh-general}
  y_h(x)=\sum_{k=1}^{r}\bigl(c_{k,0}+c_{k,1}\,x+\dots+c_{k,m_k-1}\,x^{m_k-1}\bigr)\,\eul^{\lambda_kx}.
\end{equation}
\end{teorema}

\begin{proof}
Por la conmutatividad (\cref{prop:factorizacion}) podemos escribir el
factor $(D-\lambda_k)^{m_k}$ en último lugar, de modo que actúe primero:
$\Ker(D-\lambda_k)^{m_k}\subseteq\Ker\Lop$. Por la \cref{prop:nucleo} esto
aporta $m_k$ soluciones, $\sum_km_k=n$ en total, independientes por el
\cref{lem:independencia}. Como $\dim\Ker\Lop=n$ (\cref{teo:dimension}),
forman una base.
\end{proof}

\begin{observacion}[Raíces próximas \cite{TEMARIO}]\label{obs:proximas}
El factor $x$ aparece de forma natural como límite. Si $\lambda_1\neq\lambda_2$
son raíces, el cociente
$\bigl(\eul^{\lambda_2x}-\eul^{\lambda_1x}\bigr)/(\lambda_2-\lambda_1)$ es
solución; cuando $\lambda_2\to\lambda_1$ tiende a
$\partial_\lambda\eul^{\lambda x}\big|_{\lambda_1}=x\,\eul^{\lambda_1x}$, que es la
segunda solución asociada a la raíz doble.
\end{observacion}

\begin{ejemplo}[$y'''-4y''+5y'-2y=0$ \cite{TEMARIO}]\label{ej:triple}
$p(\lambda)=\lambda^{3}-4\lambda^{2}+5\lambda-2=(\lambda-1)^{2}(\lambda-2)$: raíz
$1$ doble y raíz $2$ simple. Por el \cref{teo:base},
$y=(k_1+k_2x)\,\eul^{x}+k_3\,\eul^{2x}$.
\end{ejemplo}

\begin{ejemplo}[Vuelta al caso crítico]\label{ej:critico}
Para $\nu=2$ en el \cref{ej:amortiguado}, $p(\lambda)=(\lambda+\omega)^{2}$ y
\begin{equation}\label{eq:critico}
  y=(c_0+c_1x)\,\eul^{-\omega x},
\end{equation}
que es el regreso al equilibrio más rápido sin oscilar (\cref{fig:amortiguado}\,(b)).
\end{ejemplo}

\begin{observacion}[Raíces complejas múltiples]\label{obs:complejas-multiples}
Combinando el \cref{teo:base} con la \cref{prop:reales}: un par
$\alpha\pm\iu\beta$ de multiplicidad $m$ aporta las $2m$ soluciones reales
$x^{j}\eul^{\alpha x}\cos\beta x$ y $x^{j}\eul^{\alpha x}\sin\beta x$,
$0\le j<m$.
\end{observacion}
